%%%%%%%%%%%%%%%%%%%%%%%%%%%%%%%%%%%%%%%%%%%%%%%%%%%%%%%%%%%%%%%%%%%%%%%%%%%%%%%%%%%%%%%%%
% Autor:        Aguilar Enriquez, Paul Sebastian a.k.a. Penserbjorne
% Fecha:        05/02/2017
% Descripcion:  Plantilla base para actividades o tareas.
%%%%%%%%%%%%%%%%%%%%%%%%%%%%%%%%%%%%%%%%%%%%%%%%%%%%%%%%%%%%%%%%%%%%%%%%%%%%%%%%%%%%%%%%%

\documentclass[a4paper,11pt]{article}                 % Papel tamaño carta, texto de 11pt.

\usepackage[top=2cm, bottom=2cm, left=2.2cm, right=2.2cm]{geometry} % Margenes
\usepackage[T1]{fontenc}                              % Indicamos la codificacion de las fuentes.
\usepackage[utf8x]{inputenc}                          % Definimos la codificacion.
\usepackage{lmodern}                                  % Para poder usar acentos.
\usepackage[spanish]{babel}                           % Usaremos idioma español.
\usepackage{amsmath}                                  % Para formulas matematicas.
\usepackage{graphicx}                                 % Para imagenes.
\usepackage{float}                                    % Para posicionar objetos.
\usepackage{booktabs}                                 % Para formatear tablas.
\usepackage{hyperref}                                 % Para enlaces y referencias.

%%%%%%%%%%%%%%%%%%%%%%%%%%%%%%%%%%%%%%%%%%%%%%%%%%%%%%%%%%%%%%%%%%%%%%%%%%%%%%%%%%%%%%%%%

% Los logos tienen posiciones relativas al nombre de la escuela.
% Cada imagen esta desplazada con respecto al texto, en este caso nombre de la univseridad.
% No se necesitan paquetes adicionales, el entorno estandar para imagenes de LaTeX puede hacerlo.
% El truco esta en definir una imagen de tamaño cero, asi no afecta al centrar los titulos.
\def\logoUNAM{%
  \begin{picture}(0,0)\unitlength=1cm
    \put (-3.5,-3) {\includegraphics[width=8em]{images/escudo-unam}}
  \end{picture}
}

\def\logoFI{%
  \begin{picture}(0,0)\unitlength=1cm
    \put (0.5,-3) {\includegraphics[width=8em]{images/escudo-fi}}
  \end{picture}
}

%%%%%%%%%%%%%%%%%%%%%%%%%%%%%%%%%%%%%%%%%%%%%%%%%%%%%%%%%%%%%%%%%%%%%%%%%%%%%%%%%%%%%%%%%

%%%%%%%\author{Apellido 1 Apellido 2, Nombre1 Nombre2 --- Cuenta}  % Autor de la actividad.
\title{Lista de actividades propuestas \\ para la Feria de Agrupaciones Estudiantiles}                % Titulo de la actividad.
\date{12/02/2018}                                           % Fecha de entrega.
\def\universidad{Universidad Nacional Autónoma de México}   % Nombre de la universidad.
\def\facultad{Facultad de Ingeniería}                              % Nombre de la facultdad.
\def\semestre{2018-2}                                     % Semestre lectivo.
\def\materia{Laboratorio de Investigación y Desarrollo del Software Libre}               % Nombre de la materia y grupo.
\makeatletter

%%%%%%%%%%%%%%%%%%%%%%%%%%%%%%%%%%%%%%%%%%%%%%%%%%%%%%%%%%%%%%%%%%%%%%%%%%%%%%%%%%%%%%%%%

\begin{document}
  
  % Titulo del documento con logos.
  \begin{center}
    \logoUNAM {\Large \universidad} \logoFI\par
    {\large \facultad}\par

    \materia\par
    \semestre\par
    \@author\par
    \@date\par
    \@title
  \end{center}

  \hrulefill\par

  \pagenumbering{gobble}                              % Oculta el numero de pagina.
  \tableofcontents                                    % Crea el indice o tabla de contenido.

%%%%%%%%%%%%%%%%%%%%%%%%%%%%%%%%%%%%%%%%%%%%%%%%%%%%%%%%%%%%%%%%%%%%%%%%%%%%%%%%%%%%%%%%%

 % \newpage                                            % Inserta una pagina nueva.
  \pagenumbering{arabic}                              % Muestra el numero de pagina.
  
  \section{Proyección de OpenMovies.}                                     % Insertamos nueva seccion, SI aparece en la tabla de contenido.
  
  Esta actividad consiste en la proyección de OpenMovies durante la feria, antes y después de cada proyección se explicará cuál es la filosofía detrás de este tipo de películas, qué herramientas se utilizan y su proceso de producción.
  \paragraph{}
 \textbf{Elementos necesarios para llevar a cabo esta actividad.}
  \begin{itemize}
    \item Cañon.
    \item Bocinas.
    \item Espacio para proyección.
    \item Asientos para los asistentes.
  \end{itemize}
  
  \textbf{Duración.}
  \begin{itemize}
    \item Las películas duran entre 15 y 45 minutos.
    \item Se esperan proyectar de 3 a 5 películas diarias*. *Se trabaja aún en el listado.
  \end{itemize}
  
  
    \section{Installfest.}                                     % Insertamos nueva seccion, SI aparece en la tabla de contenido.
    Esta actividad consiste en promover el uso e instalación de distribuciones GNU Linux para uso personal y académico. Se asesorará de acuerdo a las necesidades de cada persona cuál es la distribución que más se adecua e ella.
  \paragraph{}
   \textbf{Elementos necesarios para llevar a cabo esta actividad.}
  \begin{itemize}
    \item Conexión a internet.
    \item Puntos de conexión a la red eléctrica (tomacorrientes).
    \item Mesas.
    \item USB's de 2 GB a 4 GB.
  \end{itemize}
  
  \textbf{Duración.}
  \begin{itemize}
    \item Todo el día, los tres días. Tentativo 12:00 - 20:00 hrs.
  \end{itemize}
  
      \section{Taller de OpenScad.}                                     % Insertamos nueva seccion, SI aparece en la tabla de contenido.

  Este taller consiste en la presentación de OpenScad para el modelado parametrizable 2D y 3D, el cuál sirve para la elaboración de planos que puedan ser manufacturados en máquinas de diseño (cortadora láser e impresora 3D).
  Temario: \url{https://lidsol.net/talleres/0003_openscad_basico.html}
      \paragraph{}
  \textbf{Elementos necesarios para llevar a cabo esta actividad.}
  \begin{itemize}
    \item Conexión a internet.
    \item Puntos de conexión a la red eléctrica (tomacorrientes).
    \item Espacio para el taller (mesas y sillas).
  \end{itemize}
  
  \textbf{Duración.}
  \begin{itemize}
    \item Considerando abarcar sólo los primeros tres temas; 3 horas, 3 días seguidos.
  \end{itemize}
  
            \textbf{Ponente.}
  \begin{itemize}
    \item Pablo Vivar.
  \end{itemize}
  
  
        \section{Taller de Vagrant.}                                     % Insertamos nueva seccion, SI aparece en la tabla de contenido.

  Este taller es para mostrar como crear entornos virtuales para desarrollo de software (containers).
      \paragraph{}
  \textbf{Elementos necesarios para llevar a cabo esta actividad.}
  \begin{itemize}
    \item Conexión a internet.
    \item Puntos de conexión a la red eléctrica (tomacorrientes).
    \item Espacio para el taller (mesas y sillas).
  \end{itemize}
  
  \textbf{Duración.}
  \begin{itemize}
    \item Dos horas. (Falta confirmar)
  \end{itemize}
  
              \textbf{Ponente.}
  \begin{itemize}
    \item Andrés Hernández.
  \end{itemize}
  
  

  
  
            \section{Taller de monitoreo y administración de una impresora 3D.}                                     % Insertamos nueva seccion, SI aparece en la tabla de contenido.

   En esta actividad se pretende mostrar todo lo que se necesita para monitorear y administrar una impresora 3D.
      \paragraph{}
  \textbf{Elementos necesarios para llevar a cabo esta actividad.}
  \begin{itemize}
    \item Conexión a internet.
    \item Espacio para el taller (mesas y sillas).
  \end{itemize}
  
  \textbf{Duración.}
  \begin{itemize}
    \item Tres horas. (Falta confirmar)
  \end{itemize}
  
              \textbf{Ponente.}
  \begin{itemize}
    \item Emilio Cabrera.
  \end{itemize}
  
  
              \section{Taller de KiCad.}                                     % Insertamos nueva seccion, SI aparece en la tabla de contenido.

   Este taller consiste en la presentación de KiCad para hacer placas PCB de circuitos electrónicos.
      \paragraph{}
  \textbf{Elementos necesarios para llevar a cabo esta actividad.}
  \begin{itemize}
    \item Conexión a internet.
    \item Espacio para el taller (mesas y sillas).
  \end{itemize}
  
  \textbf{Duración.}
  \begin{itemize}
    \item Dos horas, los tres días.
  \end{itemize}
  
              \textbf{Ponente.}
  \begin{itemize}
    \item Yesica Navarro.
  \end{itemize}
  
  
                \section{Taller de Nightly.}                                     % Insertamos nueva seccion, SI aparece en la tabla de contenido.

   En este taller se mostrará como instalar y configurar Firefox Nightly, se explicará la importancia de contribuir con pruebas en un software en etapa beta, cómo probarlo y reportar bugs. 
      \paragraph{}
  \textbf{Elementos necesarios para llevar a cabo esta actividad.}
  \begin{itemize}
    \item Conexión a internet.
    \item Espacio para el taller (mesas y sillas).
  \end{itemize}
  
  \textbf{Duración.}
  \begin{itemize}
    \item Dos horas, los tres días.
  \end{itemize}
  
              \textbf{Ponente.}
  \begin{itemize}
    \item Paul Aguilar.
  \end{itemize}
  
  \vspace{1 cm}
  
            \section{Conferencia/Plática sobre Derechos Digitales. (Falta confirmar)}                                     % Insertamos nueva seccion, SI aparece en la tabla de contenido.

   En esta conferencia se abordará el contexto actual de los derechos digitales, los riesgos y amenazas que existen en torno a ellos.
      \paragraph{}
  \textbf{Elementos necesarios para llevar a cabo esta actividad.}
  \begin{itemize}
    \item Auditorio.
    \item Conexión a internet. (Transmisión en vivo)
  \end{itemize}
  
  \textbf{Duración.}
  \begin{itemize}
    \item Dos horas. (Falta confirmar)
  \end{itemize}
  
    \textbf{Ponente.}
  \begin{itemize}
    \item Jacobo Nájera.
  \end{itemize}
  
           \section{Conferencia/Plática sobre Redes de Anonimato. (Falta confirmar)}                                     % Insertamos nueva seccion, SI aparece en la tabla de contenido.

   En esta plática se pretende hablar sobre mecánismos de privacidad y anonimato en la red.
      \paragraph{}
  \textbf{Elementos necesarios para llevar a cabo esta actividad.}
  \begin{itemize}
    \item Auditorio.
        \item Conexión a internet. (Transmisión en vivo)
  \end{itemize}
  
  \textbf{Duración.}
  \begin{itemize}
    \item (Falta confirmar)
  \end{itemize}
  
      \textbf{Ponente.}
  \begin{itemize}
    \item Gunnar Wolf.
  \end{itemize}
  
  \section{Conferencia/Plática de GitHub.}                                     % Insertamos nueva seccion, SI aparece en la tabla de contenido.

   En esta plática se hablará de la importancia de utilizar un control de versiones para proyectos universitarios.
      \paragraph{}
  \textbf{Elementos necesarios para llevar a cabo esta actividad.}
  \begin{itemize}
    \item Auditorio.
        \item Conexión a internet. (Transmisión en vivo)
  \end{itemize}
  
  \textbf{Duración.}
  \begin{itemize}
    \item Dos horas. (Falta confirmar)
  \end{itemize}
  
        \textbf{Ponente.}
  \begin{itemize}
    \item Pablo Flores.
  \end{itemize}
  
    \vspace{1 cm}
  
    \section{Conferencia/Plática de Mozilla.(Falta confirmar)}                                     % Insertamos nueva seccion, SI aparece en la tabla de contenido.

   Pendiente.
      \paragraph{}
  \textbf{Elementos necesarios para llevar a cabo esta actividad.}
  \begin{itemize}
    \item Auditorio.
        \item Conexión a internet. (Transmisión en vivo)
  \end{itemize}
  
  \textbf{Duración.}
  \begin{itemize}
    \item Dos horas. (Falta confirmar)
  \end{itemize}
  
          \textbf{Ponente.}
  \begin{itemize}
    \item (Falta confirmar)
  \end{itemize}
  
      \section{Conferencia/Plática de Software Libre.}                                     % Insertamos nueva seccion, SI aparece en la tabla de contenido.

   En esta plática se pretende hablar de qué es el FOSS y su impacto en el mundo. 
      \paragraph{}
  \textbf{Elementos necesarios para llevar a cabo esta actividad.}
  \begin{itemize}
    \item Auditorio.
        \item Conexión a internet. (Transmisión en vivo)
  \end{itemize}
  
  \textbf{Duración.}
  \begin{itemize}
    \item Dos horas. (Falta confirmar)
  \end{itemize}
  
            \textbf{Ponentes.}
  \begin{itemize}
    \item Miembros de LIDSOL.
  \end{itemize}
%%%%%%%%%%%%%%%%%%%%%%%%%%%%%%%%%%%%%%%%%%%%%%%%%%%%%%%%%%%%%%%%%%%%%%%%%%%%%%%%%%%%%%%%%
  


%%%%%%%%%%%%%%%%%%%%%%%%%%%%%%%%%%%%%%%%%%%%%%%%%%%%%%%%%%%%%%%%%%%%%%%%%%%%%%%%%%%%%%%%%

\end{document}
